\documentclass[12pt,a4paper]{article}

% ------------------------------------------------
% Pacotes
% ------------------------------------------------
\usepackage[utf8]{inputenc}
\usepackage[T1]{fontenc}
\usepackage[brazil]{babel}
\usepackage{graphicx}
\usepackage[table]{xcolor}
\usepackage{array}
\usepackage{tabularx}
\usepackage{geometry}
\usepackage{fancyhdr}
\usepackage{mathptmx} % Times New Roman
\usepackage{enumitem}\usepackage{amsfonts}
\usepackage{amsmath} 
\usepackage{amssymb}
\usepackage{tikz}
\usetikzlibrary{positioning, shapes.geometric}
\usepackage{listings}
\usepackage[most]{tcolorbox}
\definecolor{maincolor}{RGB}{0,100,50}
\definecolor{forestgreen}{RGB}{0,100,50}
\usepackage{booktabs}
\usepackage{pdflscape}
\usepackage{multicol}

\newtcolorbox{codeblock}{%
  colback=black!2,
  colframe=gray!30,
  arc=1mm,
  boxrule=0.5pt,
  left=5pt,
  right=5pt,
  top=5pt,
  bottom=5pt,
  fontupper=\small
}

\lstset{
    language=Java,
    basicstyle=\small\ttfamily,
    keywordstyle=\bfseries,
    showstringspaces=false,
    breaklines=true,
    frame=none,
    numbers=none,
    tabsize=2,
    commentstyle=\itshape\color{gray},
    morekeywords={this}
}

% TikZ Styles for ER Diagrams
\tikzset{
    block/.style={rectangle, draw, thick, minimum width=2cm, minimum height=1cm, fill=white},
    diamond_node/.style={diamond, draw, thick, minimum width=2cm, minimum height=1cm, fill=white, aspect=2},
    attribute/.style={ellipse, draw, thin, minimum width=1.5cm, minimum height=0.6cm, fill=white, font=\scriptsize},
    line/.style={draw, thick}, 
}
% ------------------------------------------------ 
% Configuração de página
% ------------------------------------------------
\geometry{
    a4paper,
    left=2cm,
    right=2cm,
    top=2cm,
    bottom=2cm,
    headheight=15pt 
}
% Caixa para resposta discursiva
\newcommand{\resposta}[1]{
\noindent
\framebox[\textwidth]{
\begin{minipage}{0.98\textwidth}
\vspace{#1}
\end{minipage}
}
}

\pagestyle{fancy}
\fancyhf{}
\renewcommand{\headrulewidth}{0pt}

\setlength{\parindent}{0pt}

% ------------------------------------------------
% Comandos auxiliares
% ------------------------------------------------
\newcolumntype{L}[1]{>{\raggedright\arraybackslash}p{#1}}
\newcolumntype{C}[1]{>{\centering\arraybackslash}p{#1}}

% ------------------------------------------------
% Documento
% ------------------------------------------------
\begin{document}

% ------------------------------------------------
% Cabeçalho
% ------------------------------------------------
\begin{tabularx}{\textwidth}{@{} L{4cm} X @{}}
    \raisebox{-0.5cm}{\includegraphics[width=4cm]{Universidade_de_Rio_Verde.jpg}} &
    \hfill
    \begin{tabular}{|>{\centering\arraybackslash}m{10cm}|>{\raggedright\arraybackslash}p{1.8cm}|}
        \hline
        \cellcolor{black}\color{white}\rule{0pt}{1.2cm}\textbf{AVALIAÇÃO 2 - 2ª CHAMADA} 
        & 
        \small Nota \tabularnewline
        \hline
    \end{tabular} 
\end{tabularx}

% ------------------------------------------------  
% Tabela principal
% ------------------------------------------------
\begin{tabularx}{\textwidth}{|X|L{2cm}|}
\hline
\textbf{Curso} & \textbf{Campus} \\
\textit{Engenharia de Software} & \textit{Rio Verde} \\
\hline
\multicolumn{2}{|l|}{ \textit{Lógica para Engenharia de Software}} \\
\hline
\textbf{Nome do(a) acadêmico(a)} & \textbf{Valor} \\
ISABELLE CRUVINEL SANTOS OLIVEIRA& 7,0 \\
\hline
\textbf{Professor} & \textbf{Turma} \\
\textit{ Sergio Souza Novak} & \textit{A} \\
\hline
\textbf{Data da Avaliação} & \textbf{Código} \\
\textit{ 18/05/2026} & \textit{ESW418} \\
\hline
\end{tabularx}
 
% ------------------------------------------------
% Barra de Atenção
% ------------------------------------------------
\vspace{0.2cm}

\begin{tabularx}{\textwidth}{|X|}
\hline
\cellcolor{black}\color{white}
\centering \textbf{ATENÇÃO:} Somente serão passíveis de \textbf{REVISÃO} avaliações resolvidas à \textbf{TINTA}. \tabularnewline
\hline
\end{tabularx}

\textbf{ORIENTAÇÕES PARA A RESOLUÇÃO}  

\begin{itemize}
    \item A avaliação é individual, com consulta apenas ao material escrito a mão (à caneta);
    \item A interpretação de texto faz parte da avaliação;
\end{itemize}

\begin{center}
    \rule{\textwidth}{0.4pt}
    \textbf{QUESTÕES}
\end{center}

\vspace{0.5cm} 

\begin{enumerate}

    \item Analise o diagrama de conjuntos abaixo e responda às questões:
    
    \begin{center}
    \begin{minipage}[t]{0.45\textwidth}
        \vspace{0pt}
        \centering
        \begin{tikzpicture}[>=stealth, thick, scale=1.0]
            \draw[fill=blue!10, draw=black] (0,0) ellipse (0.8 and 1.8);
            \draw[fill=red!10, draw=black] (4,0) ellipse (0.8 and 1.8);
            
            \node at (0,2.1) {\textbf{Domínio (A)}};
            \node at (4,2.1) {\textbf{Contradomínio (B)}};
    
            \node (p1) at (0, 1.2) {a};
            \node (p2) at (0, 0.4) {b};
            \node (p3) at (0, -0.4) {c};
            \node (p4) at (0, -1.2) {d};
            
            \node (q1) at (4, 1.2) {10};
            \node (q2) at (4, 0.4) {20};
            \node (q3) at (4, -0.4) {30};
            \node (q4) at (4, -1.2) {40};
    
            \draw[->] (p1) -- (q1);
            \draw[->] (p2) -- (q2);
            \draw[->] (p3) -- (q2);
            \draw[->] (p4) -- (q3);
        \end{tikzpicture}
    \end{minipage}%
    \hfill
        \vspace{1cm}
        Com base no diagrama:
        \begin{enumerate}[label=\alph*)]
            \item Qual é o conjunto \textbf{Imagem} desta relação? \\ \resposta{3cm}
            \item Esta relação é uma \textbf{função}? Justifique. \\ \resposta{3cm}
            \item Se for função, ela é \textbf{injetora}? Por quê? \\ \resposta{3cm}
            \item Se for função, ela é \textbf{sobrejetora}? Por quê? \\ \resposta{3cm}
        \end{enumerate}
    \end{center}

    \item Em um sistema de \textbf{e-commerce}, a finalização de um pedido ocorre se certas condições lógicas forem atendidas.
    As condições são: \texttt{inStock} (produto em estoque), \texttt{paymentOk} (pagamento aprovado), \texttt{isExpress} (entrega expressa) e \texttt{isVip} (cliente VIP).
    
    A lógica de validação é definida por:
    
    \begin{codeblock}
    \begin{verbatim}
    bool canProcessOrder() {
        return ( (inStock \&\& paymentOk) || isVip ) &&
               ( (inStock || isVip) && (paymentOk || isVip) );
    }
    \end{verbatim}
    \end{codeblock}

    Utilizando as propriedades da Álgebra Booleana, assinale a alternativa que apresenta a \textbf{lógica mínima simplificada} equivalente à expressão acima:

    \begin{enumerate}[label=\Alph*)]
        \item \texttt{(inStock \&\& paymentOk) || isVip}
        \item \texttt{inStock \&\& (paymentOk || isVip)}
        \item \texttt{(inStock || paymentOk) \&\& isVip}
        \item \texttt{inStock \&\& paymentOk \&\& isVip}
        \item \texttt{inStock || paymentOk || isVip}
    \end{enumerate}

    \item Sobre a classificação de funções matemáticas, analise as afirmações abaixo:
    
    \begin{enumerate}[label=\Roman*.]
        \item Uma função é bijetora quando é simultaneamente injetora e sobrejetora.
        \item Em uma função sobrejetora, o conjunto imagem é igual ao contradomínio.
        \item Se uma função é injetora, dois elementos distintos do domínio podem ter a mesma imagem.
        \item Uma função constante $f(x) = k$ é sempre injetora para qualquer domínio com mais de um elemento.
    \end{enumerate}

    Assinale a alternativa que indica quais afirmações estão corretas:

    \begin{enumerate}[label=\Alph*)]
        \item Apenas as afirmações I e II estão corretas.
        \item Apenas as afirmações I e III estão corretas.
        \item Apenas as afirmações II e IV estão corretas.
        \item Apenas as afirmações III e IV estão corretas.
        \item Todas as afirmações estão corretas.
    \end{enumerate}

    \item Um sistema de \textbf{Irrigação Inteligente} para plantações de soja utiliza sensores IoT para decidir o acionamento das bombas. Os sensores são:
    
    \begin{itemize}[leftmargin=1cm, label=\textbullet]
      \item \textbf{A} — Sensor de Umidade do Solo: $A = 1$ indica solo seco.
      \item \textbf{B} — Sensor de Nível do Reservatório: $B = 1$ indica reservatório com água suficiente.
      \item \textbf{C} — Relógio de Controle: $C = 1$ indica que está no horário permitido para irrigação.
      \item \textbf{D} — Sensor de Chuva: $D = 1$ indica que está chovendo no momento.
    \end{itemize}
    
    \textbf{Saída do sistema:} $Y = 1$ liga a irrigação; $Y = 0$ mantém desligada.
    
    \noindent\textbf{Tabela Verdade:}
    
    \medskip
    \begin{minipage}[c]{0.5\textwidth}
        \centering
        \renewcommand{\arraystretch}{1.2}
        \begin{tabular}{|c|c|c|c||c|}
          \hline
          \textbf{A} & \textbf{B} & \textbf{C} & \textbf{D} & \textbf{Y} \\
          \hline
          0 & 0 & 0 & 0 & 0 \\
          0 & 0 & 0 & 1 & 0 \\
          0 & 1 & 0 & 0 & 0 \\
          0 & 1 & 1 & 0 & 0 \\
          1 & 0 & 0 & 0 & 0 \\
          1 & 1 & 0 & 0 & 0 \\
          1 & 1 & 0 & 1 & 0 \\
          1 & 1 & 1 & 0 & 1 \\
          1 & 1 & 1 & 1 & 0 \\
          0 & 0 & 1 & 1 & 0 \\
          0 & 1 & 0 & 1 & 0 \\
          1 & 0 & 1 & 0 & 1 \\
          1 & 0 & 1 & 1 & 1 \\
          1 & 1 & 1 & 1 & 1 \\
          \multicolumn{5}{|c|}{\textit{(Linhas omitidas têm $Y=0$ por padrão)}} \\
          \hline
        \end{tabular}
    \end{minipage}
    \hfill
    \begin{minipage}[c]{0.45\textwidth}
        \textbf{Tarefas:}\\
        \begin{enumerate}[label=\textbf{\alph*)}, leftmargin=1cm]
          \item Escreva a função booleana baseada na tabela acima.
          \item Simplifique a expressão utilizando Álgebra Booleana.
          \item Desenhe o circuito lógico correspondente.
        \end{enumerate}
    \end{minipage}
    
    \resposta{12cm}

    \item Um sistema de backup precisa realizar 4 tarefas diferentes em uma ordem específica. De quantas maneiras distintas o sistema pode organizar o agendamento dessas 4 tarefas?
    
    \begin{enumerate}[label=\Alph*)]
        \item 4
        \item 12
        \item 16
        \item 24
    \end{enumerate}
    \item Analise os gráficos abaixo e marque \textbf{V} para verdadeiro ou \textbf{F} para falso para as classificações de cada um deles. Considere o domínio e o contradomínio como $\mathbb{R}$:

    \begin{center}
    \renewcommand{\arraystretch}{1.5}
    \begin{tabular}{cc}
        \begin{tikzpicture}[scale=0.8]
            \draw[->] (-2,0) -- (2,0) node[right] {$x$};
            \draw[->] (0,-1) -- (0,3) node[above] {$y$};
            \draw[thick, red, domain=-1.8:1.8] plot (\x, {abs(\x)});
        \end{tikzpicture} & 
        \begin{tikzpicture}[scale=0.8]
            \draw[->] (-2,0) -- (2,0) node[right] {$x$};
            \draw[->] (0,-2) -- (0,2) node[above] {$y$};
            \draw[thick, red, domain=-1.4:1.4] plot (\x, {0.5*\x*\x*\x});
        \end{tikzpicture} \\
        \small I. \begin{tabular}{ll} (\quad) & Não é função \\ (\quad) & Função Injetora \\ (\quad) & Função Sobrejetora \end{tabular} &
        \small II. \begin{tabular}{ll} (\quad) & Não é função \\ (\quad) & Função Injetora \\ (\quad) & Função Sobrejetora \end{tabular} \\ [1.5em]
        \begin{tikzpicture}[scale=0.8]
            \draw[->] (-2,0) -- (2,0) node[right] {$x$};
            \draw[->] (0,-2) -- (0,2) node[above] {$y$};
            \draw[thick, red, smooth, samples=100, domain=0:2*pi] plot ({1.2*sin(\x r)^3}, {1.1*cos(\x r) - 0.4*cos(2*\x r) - 0.15*cos(3*\x r) - 0.05*cos(4*\x r)});
        \end{tikzpicture} & 
        \begin{tikzpicture}[scale=0.8]
            \draw[->] (-2,0) -- (2,0) node[right] {$x$};
            \draw[->] (0,-0.5) -- (0,2.5) node[above] {$y$};
            \draw[thick, red, domain=-2:1.2] plot (\x, {exp(\x)});
        \end{tikzpicture} \\
        \small III. \begin{tabular}{ll} (\quad) & Não é função \\ (\quad) & Função Injetora \\ (\quad) & Função Sobrejetora \end{tabular} &
        \small IV. \begin{tabular}{ll} (\quad) & Não é função \\ (\quad) & Função Injetora \\ (\quad) & Função Sobrejetora \end{tabular}
    \end{tabular}
    \end{center}

    \item Desenvolva a expansão do seguinte binômio utilizando o \textbf{Binômio de Newton}:
    \[ \left( \sin(x) - \frac{1}{\sqrt{x}}  \right)^6 \]
    Apresente todos os termos e simplifique os coeficientes.
    
    \resposta{6cm}

    \item Determine a \textbf{soma dos coeficientes} do desenvolvimento de $(2x + y)^6$.
    
    \begin{enumerate}[label=\Alph*)]
        \item 64
        \item 243
        \item 729
        \item 4096
    \end{enumerate}



\end{enumerate}

\newpage

% ------------------------------------------------
% ANEXO: FORMULÁRIO INTEGRADO
% ------------------------------------------------
\begin{center}
  {\Large \textbf{Formulário de Apoio}}
\end{center}

\vspace{-0.2cm}

\noindent \textbf{1. Propriedades da Álgebra Booleana}
\begin{center}
\renewcommand{\arraystretch}{1.1}
\footnotesize
\begin{tabular}{
  >{\centering\arraybackslash}m{0.4cm}
  >{\centering\arraybackslash}m{4.3cm}
  >{\centering\arraybackslash}m{4.3cm}
  >{\centering\arraybackslash}m{4.5cm}
}
\toprule
\textbf{N°} & \textbf{Forma AND (·)} & \textbf{Forma OR (+)} & \textbf{Nome da Propriedade} \\
\midrule
1 & $A \cdot 0 = 0$ & $A + 1 = 1$ & Elemento Nulo \\
\midrule
2 & $A \cdot 1 = A$ & $A + 0 = A$ & Elemento Neutro (Identidade) \\
\midrule
3 & $A \cdot A = A$ & $A + A = A$ & Idempotência \\
\midrule
4 & $A \cdot \overline{A} = 0$ & $A + \overline{A} = 1$ & Complemento \\
\midrule
5 & $A \cdot B = B \cdot A$ & $A + B = B + A$ & Comutatividade \\
\midrule
6 & $A \cdot (B \cdot C) = (A \cdot B) \cdot C$ & $A + (B + C) = (A + B) + C$ & Associatividade \\
\midrule
7 & $\overline{A \cdot B \cdot C} = \overline{A} + \overline{B} + \overline{C}$ & $\overline{A + B + C} = \overline{A} \cdot \overline{B} \cdot \overline{C}$ & Teorema de De Morgan \\
\midrule
8 & $A \cdot (B + C) = A \cdot B + A \cdot C$ & $(A+B)\cdot(A+C) = A + B \cdot C$ & Distributividade \\
\midrule
9 & $A + A \cdot B = A$ & $A \cdot (A + B) = A$ & Absorção \\
\midrule
10 & $A + \overline{A} \cdot B = A + B$ & $A \cdot (\overline{A} + B) = A \cdot B$ & Cobertura (Absorção Complementar) \\
\midrule
11 & $\overline{\overline{A}} = A$ & $\overline{\overline{A}} = A$ & Dupla Negação (Involução) \\
\midrule
12 & $A \cdot B + \overline{A} \cdot C + B \cdot C = A \cdot B + \overline{A} \cdot C$ & $(A+B) \cdot (\overline{A}+C) \cdot (B+C) = (A+B) \cdot (\overline{A}+C)$ & Teorema do Consenso \\
\midrule
13 & $C \cdot A + C \cdot \overline{A} \cdot B = C \cdot A + C \cdot B$ & $(C+A)(C+\overline{A}+B) = (C+A)(C+B)$ & Simplificação por Complemento \\
\bottomrule
\end{tabular}
\end{center}

\begin{multicols}{2}
    \noindent \textbf{2. Análise Combinatória e Binômio}

    
    \textbf{Binômio:} $(a+b)^n = \sum_{k=0}^n \binom{n}{k} a^{n-k} b^k$ \\[0.15cm]
    \textbf{Termo Geral:} $T_{k+1} = \binom{n}{k} a^{n-k} b^k$ \\[0.15cm]
    \textbf{Permutação Simples:} $P_n = n!$ \\[0.15cm]
    \textbf{Permutação (Rep.):} $P_n^{k_1, \dots, k_m} = \frac{n!}{k_1! \dots k_m!}$ \\[0.15cm]
    \textbf{Arranjo Simples:} $A_{n,k} = \frac{n!}{(n-k)!}$ \\[0.15cm]
    \textbf{Arranjo (Rep.):} $AR_{n,k} = n^k$ \\[0.15cm]
    \textbf{Combinação Simples:} $C_{n,k} = \binom{n}{k} = \frac{n!}{k!(n-k)!}$ \\[0.15cm]
    \textbf{Combinação (Rep.):} $CR_{n,k} = \binom{n+k-1}{k}$ \\[0.15cm]
    \textbf{Multinômio:} $\frac{n!}{k_1! \dots k_m!} x_1^{k_1} \dots x_m^{k_m}$

    \vspace{0.3cm}
    \noindent \textbf{4. Funções e Relações}\\
    
    Relação ($R: A \to B$): $R \subseteq A \times B$ \\[0.1cm]
    Função ($f: A \to B$): $\forall x \in A, \exists! y \in B \mid (x, y) \in f$ \\[0.1cm]
    Injetora:$\forall x_1, x_2 \in A, f(x_1) = f(x_2) \implies x_1 = x_2$ \\[0.1cm]
    Sobrejetora: $\forall y \in B, \exists x \in A \mid f(x) = y$ \\[0.1cm]
    Bijetora: $\forall x_1, x_2 \in A, f(x_1) = f(x_2) \implies x_1 = x_2$ $\land$ $\forall y \in B, \exists x \in A \mid f(x) = y$

    \columnbreak

    \noindent \textbf{3. Triângulo de Pascal}
    \begin{center}
        \begin{tikzpicture}[x=0.42cm, y=0.38cm, font=\small,scale=1.9]
        \node at (0,0) {1};
        \node at (-0.5,-1) {1}; \node at (0.5,-1) {1};
        \node at (-1,-2) {1}; \node at (0,-2) {2}; \node at (1,-2) {1};
        \node at (-1.5,-3) {1}; \node at (-0.5,-3) {3}; \node at (0.5,-3) {3}; \node at (1.5,-3) {1};
        \node at (-2,-4) {1}; \node at (-1,-4) {4}; \node at (0,-4) {6}; \node at (1,-4) {4}; \node at (2,-4) {1};
        \node at (-2.5,-5) {1}; \node at (-1.5,-5) {5}; \node at (-0.5,-5) {10}; \node at (0.5,-5) {10}; \node at (1.5,-5) {5}; \node at (2.5,-5) {1};
        \node at (-3,-6) {1}; \node at (-2,-6) {6}; \node at (-1,-6) {15}; \node at (0,-6) {20}; \node at (1,-6) {15}; \node at (2,-6) {6}; \node at (3,-6) {1};
        \node at (-3.5,-7) {1}; \node at (-2.5,-7) {7}; \node at (-1.5,-7) {21}; \node at (-0.5,-7) {35}; \node at (0.5,-7) {35}; \node at (1.5,-7) {21}; \node at (2.5,-7) {7}; \node at (3.5,-7) {1};
        \node at (-4,-8) {1}; \node at (-3,-8) {8}; \node at (-2,-8) {28}; \node at (-1,-8) {56}; \node at (0,-8) {70}; \node at (1,-8) {56}; \node at (2,-8) {28}; \node at (3,-8) {8}; \node at (4,-8) {1};
        \node at (-4.5,-9) {1}; \node at (-3.5,-9) {9}; \node at (-2.5,-9) {36}; \node at (-1.5,-9) {84}; \node at (-0.5,-9) {126}; \node at (0.5,-9) {126}; \node at (1.5,-9) {84}; \node at (2.5,-9) {36}; \node at (3.5,-9) {9}; \node at (4.5,-9) {1};
        \node at (-5,-10) {1}; \node at (-4,-10) {10}; \node at (-3,-10) {45}; \node at (-2,-10) {120}; \node at (-1,-10) {210}; \node at (0,-10) {252}; \node at (1,-10) {210}; \node at (2,-10) {120}; \node at (3,-10) {45}; \node at (4,-10) {10}; \node at (5,-10) {1};
                    \end{tikzpicture}
    \end{center}
\end{multicols}
\end{document}
